\documentclass[11pt]{article}
\usepackage[margin=25mm]{geometry}
\usepackage{fontspec}
\setmainfont{Latin Modern Roman}
\usepackage{amsmath,amssymb,mathtools}
\usepackage{unicode-math}
\setmathfont{Latin Modern Math}
\usepackage{xurl,hyperref}
\hypersetup{colorlinks=true,urlcolor=blue}
\setcounter{secnumdepth}{0}
\setlength{\parindent}{0pt}
\setlength{\parskip}{5pt}
\setlength{\emergencystretch}{3em}
\newcommand{\sref}[2]{\hyperlink{source:#1:#2}{[S#2]}}
\newcommand{\cataloguescope}[1]{\paragraph{Recorded target.}#1}
\newcommand{\R}{\mathbb R}
\newcommand{\C}{\mathbb C}
\newcommand{\N}{\mathbb N}
\newcommand{\Z}{\mathbb Z}
\newcommand{\Q}{\mathbb Q}
\newcommand{\F}{\mathbb F}
\newcommand{\eps}{\varepsilon}
\begin{document}
% UnsolvedMath ID 30000362; source catalogue code OWR-1114-003
\setcounter{section}{30000361}
\section{Deformations and Massless Spectra of Smooth Deligne–Mumford Stacks}\label{30000362}
{\small UnsolvedMath ID 30000362 · Mathematical Physics}
\subsection{Definitions and mathematical statement}

\paragraph{Shared notation.}$\mathbb N=\{1,2,\ldots\}$; the standard number fields have their usual meanings. All additional conventions are stated locally.


\paragraph{Definitions and conventions.} A smooth Deligne--Mumford stack is a geometric object locally presented by smooth spaces with finite stabilizer groups. A quotient presentation $[X/G]$ specifies a space and a group action; distinct presentations may describe the same stack. The report associates a gauged sigma model to physically realizable presentations. A deformation changes the background or couplings infinitesimally; a massless spectrum consists of zero-mass physical excitations. The inertia stack records pairs of a point and an automorphism of that point and enters the proposed cohomological spectrum calculation.

\paragraph{Source scope and qualifications.} The original registry excerpt also includes infrared agreement of different quotient presentations, separately catalogued at ID 30000363. Both aspects of the excerpt are retained by this note. Ordinary stack deformation theory can fail to match the physical deformations, as the report's C^2/Z\_2 example demonstrates.

\paragraph{Bundled original excerpt.} The supplied original excerpt also asks whether different quotient presentations of the same stack flow to the same physical infrared theory. That presentation-independence question is recorded separately as problem 30000363; no claim that it follows from a spectrum prescription is made.

\paragraph{Catalogue mathematical target.}
\paragraph{Catalogue wording retained for traceability.}
For string compactification on a smooth Deligne–Mumford stack, how should deformations and massless spectra be defined so that they agree with the physical theory?

\paragraph{Saved literature context.} The supplied catalogue's assessment dated 2026-08-21 records: Stack presentations have been incorporated into gauged linear sigma models and gerbe/string constructions, producing concrete conformal theories in many examples; a general classification of strings and conformal field theories on stacks remains open. The assessment is retained as dated provenance; it is not a new comprehensive literature or open-status certification.
\paragraph{Scope and source qualification.}The original registry excerpt also includes infrared agreement of different quotient presentations, separately catalogued at ID 30000363. Both aspects of the excerpt are retained by this note. Ordinary stack deformation theory can fail to match the physical deformations, as the report's C^2/Z_2 example demonstrates.
\subsection{Short English statement}
Define deformations and massless spectra for string theory on stacks so that they agree with the associated physical theories.
\subsection{Research attempt}
\paragraph{Model, date and investigated scope.}
OpenAI GPT-6 Astra (configured model \texttt{gpt-6-astra}), ultra reasoning; generated 11 October 2026. Computes the inertia components and conditional age grading of the C²/Z₂ example; no physical spectrum/deformation theorem is claimed.

\paragraph{Sources and literature check.}
The complete supplied definition, including source qualifications and the dated literature assessment, was reviewed. The calculation below is an elementary restricted attempt motivated by [S2]; it does not claim novelty or a fresh literature-status determination.

\paragraph{Context and investigated scope.} The complete supplied definition, including source qualifications and the dated literature assessment, was reviewed. The calculation below is an elementary restricted attempt motivated by [S2]; it does not claim novelty or a fresh literature-status determination.

\paragraph{The smallest twisted-sector calculation.} Let $X=\mathbb C^2$ and let $G=\{1,g\}$ act by $g(z_1,z_2)=(-z_1,-z_2)$. The inertia stack consists of points together with stabilizer elements. Since $G$ is abelian, its components are
\[
 I[X/G]=[X/G]\ \sqcup\ [\{0\}/G].
\]
Indeed the identity fixes all of $X$, whereas solving $gz=z$ gives $2z=0$ and hence $z=0$ over $\mathbb C$. The second component is missed if one only keeps invariant differential forms on the original $X$.

\paragraph{Degree bookkeeping.} Each eigenvalue of $g$ on the tangent space at zero is $-1=e^{2\pi i(1/2)}$. The sum of the fractional exponents is one. If a proposed sector-cohomology prescription shifts real cohomological degree by twice this sum, its constant twisted-sector class sits in degree two. This statement is just the calculation within that explicitly specified prescription. It does not prove that the class is a normalizable massless state in the noncompact physical theory, nor that it integrates to a deformation.

\paragraph{Failure of the naive identification.} A deformation prescription using only the untwisted component cannot even represent this extra candidate sector. Conversely, adding inertia cohomology is not a proof of agreement with the physical spectrum: normalizability, projections, interactions, and obstruction equations must be supplied. The source's additional presentation-independence question, separately numbered 30000363, also remains open here; equality of this elementary stack invariant would not identify infrared theories. The first gap is constructing the physical state complex and a map from the candidate twisted class into its massless cohomology.

\paragraph{Verification and adversarial audit.}
Checked fixed loci and both tangent eigenvalues, noncompact normalizability caveat, and retained the bundled presentation-independence question. Model self-audit only; no independent human proof checking. Prompts A and B were applied to this canonical record and its qualified source scope. Checked fixed loci and both tangent eigenvalues, noncompact normalizability caveat, and retained the bundled presentation-independence question. Retained unresolved status and zero original-target progress; A physical state complex, normalizability criterion and comparison with candidate sector cohomology are missing.

Coordinator checked the two inertia components and age-one degree bookkeeping, with normalizability and deformation obstructions left unresolved. This additional check was model-assisted, not independent human verification.

\paragraph{Tools, process and contributions.}
Python 3 was used for file preparation and structural checks; no numerical experiment or theorem prover was used. Contribution prompt A was followed by prompt B. GPT-6 Astra with ultra reasoning drafted the research and performed its stated self-audit; any additional coordinator contribution is identified in the verification and audit. A GPT-6-based Codex coordinator assembled the manuscript and checked repository integration; its exact variant and reasoning configuration were not exposed. No human mathematical contribution, independent human proof check or formal proof certification is supplied. Complete derivations and the consulted source references are included here; private reasoning transcripts are not published.

\paragraph{Final status and first remaining gap.}
\textbf{UNRESOLVED}. A physical state complex, normalizability criterion and comparison with candidate sector cohomology are missing.

\section{Completion Estimate}
\noindent\textbf{COMPLETION: 0\%}
The calculation is an auxiliary or restricted attempt and establishes no newly verified portion of the full source target. A physical state complex, normalizability criterion and comparison with candidate sector cohomology are missing.

\subsection{Sources}
\begin{itemize}
\item[\hypertarget{source:30000362:1}{S1}] ULAM.AI and the UnsolvedMath Contributors, UnsolvedMath, record 30000362 (OWR-1114-003). Catalogue curation under CC BY 4.0; definitions and English exposition added locally.
\url{https://huggingface.co/datasets/ulamai/UnsolvedMath}.
\item[\hypertarget{source:30000362:2}{S2}] Mini-Workshop: Heterotic Strings, Derived Categories, and Stacks (2005), Oberwolfach Reports, DOI 10.4171/owr/2005/53, pp. 3019--3060; official report source and mathematical context.
\url{https://doi.org/10.4171/owr/2005/53}.
\item[\hypertarget{source:30000362:3}{S3}] Official EMS Press PDF of the same Oberwolfach report; PDF pages 18--20, printed pages 3036--3038.
\url{https://ems.press/content/serial-article-files/46027}.
\item[\hypertarget{source:30000362:4}{S4}] GLSMs for gerbes and stacks; reference in the saved dated catalogue literature review.
\url{https://arxiv.org/abs/hep-th/0502053}.
\item[\hypertarget{source:30000362:5}{S5}] derived categories and stacks in physics; reference in the saved dated catalogue literature review.
\url{https://doi.org/10.48550/arXiv.hep-th/0608056}.
\end{itemize}
\label{end:30000362}
\end{document}
